\documentclass{covermaker}
\doctitle{Software Requirements Specification}
\docsubtitle{Requirements for Cover\_Maker v1.0 (after ISO/IEC/IEEE 29148)}
\docid{CM-SRS}
\docstatus{Released}
\input{glossary}

\begin{document}
\cmtitlepage

\begin{revisionhistory}
  0.1 & 2026-11-03 & Initial requirements catalogue and constraints (M0). \\
  1.0 & 2026-11-05 & Version 1.0: DAW and discovery requirements (M11); traceability to the Test Report. \\
\end{revisionhistory}

\tableofcontents

\chapter{Introduction}

\section{Purpose}
This \gls{srs} defines what \CM{} version 1.0 shall do and the constraints it
must respect. It is the reference for design, testing and acceptance.

\section{Scope}
\CM{} is a Windows desktop application that separates audio files supplied by
the user into selectable \glspl{stem} and renders \glspl{backingtrack}, entirely
on the local machine.

\section{Conventions}
Requirement identifiers are stable: \texttt{CON} for constraints, \texttt{FR}
for functional and \texttt{NFR} for non-functional requirements. Priorities are
\emph{must}, \emph{should} and \emph{could}. The catalogue is maintained in
\texttt{docs/data/requirements.yaml}; the tables below are generated from it.

\section{References}
Project Plan (CM-PP); Architecture \& Design Document (CM-ADD).

\chapter{Overall description}

\section{Product perspective}
\CM{} consists of a desktop shell (Tauri 2), a user interface (SvelteKit) and a
local Python engine that runs the separation models on the CPU or an NVIDIA
\gls{gpu}. It has no server component; network access is limited to downloading
models chosen by the user, application updates and optional Creative Commons
catalogues.

\section{Users}
\begin{description}
  \item[Musician] Practises or records covers; needs quick backing tracks
        without specific instruments.
  \item[Producer] Works in a \gls{daw}; needs clean stems from material they
        have rights to, without disturbing a running session.
\end{description}

\section{Operating environment}
Windows 10 or 11, 64-bit; 16 GB RAM minimum; optional NVIDIA GPU (Turing or
newer) with at least 4 GB \gls{vram}; WebView2 runtime.

\section{Assumptions}
Users process material they own or are licensed to use. Model availability
depends on third-party hosts.

\chapter{Constraints}
\input{generated/constraints}

\chapter{Functional requirements}
\input{generated/requirements-functional}

\chapter{Non-functional requirements}
\input{generated/requirements-nonfunctional}

\chapter{Traceability}
The matrix below is generated from \texttt{docs/data/traceability.yaml}; an
engine test fails if a requirement has no verification or cites a test that
does not exist. The Test Report records the results.

{\small
\input{generated/traceability}
}

\printnoidxglossaries
\end{document}
